Lima (1973, \textit{apud} \cite{andreataAulaExpositivaPaulo2019}) contextualiza as origens do método expositivo, processo primordial de comunicação oral; como anterior à descoberta da escrita e da imprensa. A exemplo do quão antiga é a ferramenta, propõe uma ambientação no contexto do ensino universitário medieval, que dispunha apenas deste recurso para ensinar os alunos. Em particular, Lima(1973, \textit{apud} \cite{andreataAulaExpositivaPaulo2019}) descreve o cenário para o docente era o seguinte: ``A inexistência de livros levava-o a repetir (recitações-leitura) o texto tantas vezes quantas necessárias para que os alunos o decorassem (hipótese em que se poderia supor que os alunos podiam ser até analfabetos)''.


Primeiro pedagogo da história, o filósofo grego Platão (427-347 a.C.) fazia uso do ensino oral em suas discussões, que --- para Sócrates (469/470-399 a.C.), mestre de Platão; e os sofistas --- tinha papel não só educativo, mas de produção da própria educação; simultâneamente visando a formação do indivíduo, possível de reconhecer pela aquisição de uma ``competência discursiva''. \cite{clermont2013pedagogia} ainda complementa dizendo que, para os sofistas, ser instruído é ``saber falar'' e ``saber argumentar'' em público segundo as regras pragmáticas da retórica.

Mattos (1959 \textit{apud} \citeauthor{andreataAulaExpositivaPaulo2019}, \citeyear{andreataAulaExpositivaPaulo2019}), \citet{nerici1968introducao} e \citet{piletti1987didatica} contribuem para a definição da metodologia expositiva. Em resumo, pode-se entendê-la como técnica de exposição oral, por parte do docente, de um determinado assunto, valendo-se de uma apresentação logicamente estruturada. Nela, o aluno tem papel de ouvinte, sendo esse responsável por assimilar o conteúdo ministrado pelo docente.

Há, no entanto --- desde a Antiga Grécia --- um grande risco para essa técnica, que mais tarde será reatado por Roma e a Idade Média: o uso majoritário de discursos preparados e repetição de memória contribuiam para um saber superficial e pretencioso, como destaca Gal (1989, \textit{apud} \citeauthor{andreataAulaExpositivaPaulo2019}, \citeyear{andreataAulaExpositivaPaulo2019}).

As principais críticas à aula expositiva são centradas na passividade do aluno, que é relegado no seu próprio processo de aprendizagem. \cite{baldanRepresentacaoAtoEnsinar2011} descreve o ato de ensinar, no sentido tradicional, da seguinte forma:

\begin{customquote}
    [...] processo mecânico, no qual o professor detém o conhecimento e o aluno assume a função passiva de assimilação deste conteúdo. Os métodos são mecânicos, com repetição e cópia, aula expositiva e arguição oral, com o uso de instrumentos disciplinadores e reguladores do comportamento. \citep{baldanRepresentacaoAtoEnsinar2011}
\end{customquote}

Outros autores, como \cite{delimaCRITICAMETODOLOGIATRADICIONAL2022} acrescentam dizendo que o problema da metodologia tradiconal é a falta de interação entre o sujeito e o objeto, a falta de diálogo entre o professor e o aluno; pois muitas vezes o que é exposto não faz dimensão alguma com a realidade do aluno presente. Isso ocorre porque há uma estrutura inflexível e unilateral, em que o docente "despeja" o saber e pouco faz para adaptar essa dinâmica às peculiaridades do aluno. Tal realidade se faz mais desafiadora no contexto do Ensino à Diastância (EAD), no qual há maior distânciamento entre sujeito e objeto, compromentendo a dinâmica já comprovada ineficaz em ambientes convencionais, no caso a sala de aula \citep{vasconcellosMetodologiaDialeticaEm1992}.

Isto posto, é uma técnica já alardeada por inúmeros autores no que diz respeito à sua eficácia, inclusive na contemporâniedade, na qual há pleno aparato tecnológico para identificação de gargalos e solução dos problemas crônicos da aula expositiva. Adicionalmente, deve-se levar em conta sua popularidade atemporal; que convém o título de metodologia mais  usada em território brasileiro \citep{nerici1968introducao}.